% Generated draft: requires mathematical and source review.
\documentclass[11pt]{article}
\usepackage[margin=25mm]{geometry}
\usepackage{fontspec}
\setmainfont{Latin Modern Roman}
\usepackage{amsmath,amssymb,mathtools,unicode-math}
\setmathfont{Latin Modern Math}
\usepackage{enumitem,needspace,xurl,hyperref}
\setlength{\parindent}{0pt}
\setlength{\parskip}{5pt}
\setlength{\emergencystretch}{3em}
\newcommand{\N}{\mathbb{N}}
\newcommand{\Z}{\mathbb{Z}}
\newcommand{\Q}{\mathbb{Q}}
\newcommand{\R}{\mathbb{R}}
\newcommand{\C}{\mathbb{C}}
\newcommand{\F}{\mathbb{F}}
\newcommand{\A}{\mathbb{A}}
\newcommand{\PP}{\mathbb{P}}
\newcommand{\E}{\mathbb{E}}
\newcommand{\eps}{\varepsilon}
\newcommand{\dd}{\,\mathrm{d}}
\begin{document}
\section*{UnsolvedMath ID: 10400172}
\textbf{Source problem code:} catalogue entry
\subsection{Definitions and mathematical statement}
A fusion-rule algebra records products of finitely many simple labels with nonnegative integral structure constants. Coherent \(6j\)-symbols encode associativity data; extra normalization and duality conditions are needed for the source's state-sum/TQFT construction. Y. Kawahigashi's Problem 9.8 concerns the particular examples of such data in [14] that do not seem to arise from quantum groups, and further conjectured examples in [160]. Their defining data are not included in the record.

The main questions are:
\begin{enumerate}
\item What TQFTs correspond to those examples?
\item If the series conjectured in [160] exists, giving a parameterized family of TQFTs, does differentiation in the parameter after a suitable reparameterization produce a more interesting invariant, possibly of Vassiliev type?
\end{enumerate}
The existence of the conjectured series is conditional. A differentiation request needs a parameter domain, a meaning of analytic or formal dependence, and a specified invariant to differentiate; a discrete indexing set alone does not define a derivative. Neither these data nor the reparameterization or finite-type convention are supplied. The source's phrase ``do not seem to arise'' is not a proof excluding every quantum-group realization.

The extraction then enters Section 9.4, ``Turaev--Viro--Ocneanu invariants.'' It calls a \(3\)-manifold state sum from \(6j\)-symbols a Turaev--Viro--Ocneanu invariant when those symbols come from a subfactor. It reports infinitely many subfactors beyond those derived from quantum or finite groups and suggests their invariants may be new. N. Sato's supplied contextual discussion concerns the Haagerup subfactor, whose index is conventionally read from the flattened fraction as \((5+\sqrt{13})/2\). It describes this as the smallest index above \(4\) among finite-depth subfactors and expects exotic properties, while reporting limited distinguishing power of its manifold invariant. Specifically, it reports equality of that invariant on
\[
 L(5,1),L(5,2),\qquad\text{and on}\qquad L(7,1),L(7,2).
\]
It reports that generalized \(E_6\)-subfactors with symmetry groups \(\Z/3\Z\) and \(\Z/5\Z\) distinguish respectively
\[
 L(3,1)\text{ from }L(3,2),\qquad
 L(5,1)\text{ from }L(5,2).
\]
These are contextual source reports, not independently checked computations; lens-space orientation and invariant-normalization conventions are absent. They do not identify the full TQFTs or construct the conjectural differentiated invariant.
\subsection{Short English statement}
Identify TQFTs arising from the cited potentially exotic fusion data. If the conjectured parameterized series exists, investigate whether reparameterized differentiation gives useful invariants, possibly finite type. The following source context reports examples and limitations of subfactor-based state sums, without resolving the conditional program.
\subsection{Sources}
\begin{itemize}
\item Local UnsolvedMath registry, permanent ID \(10400172\), exact source code \texttt{catalogue entry}, Problem 9.8, attributed to Y. Kawahigashi. The following contextual remarks are attributed to N. Sato. The title is truncated after ``(Y.''
\item Ohtsuki, \emph{Problems on invariants of knots and 3-manifolds} (2002), Problem 9.8 and extracted Section 9.4 context, PDF page 140: \url{https://msp.org/gtm/2002/04/gtm-2002-04-024s.pdf}. Bold headings distinguish 231 numbered items from surrounding references.
\item Internal references [14] for the reported examples and [160] for the conjectured series; bibliography entries are absent. The Haagerup and generalized \(E_6\) computations have no additional citation in this extract.
\item Extraction pdf-font-structure, accessed 2026-07-29. Source evidence NEEDS\_REVIEW presents an open problem and calls for a current primary-literature audit. Registry status partially\_solved and classification PARTIAL-PROGRESS say the Haagerup subfactor and related constructions give exotic examples; they do not specify which further questions are settled.
\item Dataset attribution: Ulam AI / UnsolvedMath Contributors, CC BY 4.0. Source rights NEEDS\_REVIEW: publicly accessible author/publisher copy, redistribution terms requiring review.
\end{itemize}
\end{document}
